\documentclass[11pt]{article}

\usepackage[margin=1in]{geometry}
\usepackage[T1]{fontenc}
\usepackage[utf8]{inputenc}
\usepackage{lmodern}
\usepackage{array}
\usepackage{booktabs}
\usepackage{enumitem}
\usepackage{hyperref}
\usepackage{longtable}
\usepackage{xcolor}
\usepackage{listings}

\hypersetup{
  colorlinks=true,
  linkcolor=blue,
  urlcolor=blue,
  citecolor=blue
}

\definecolor{cacheirblue}{HTML}{1F4E79}
\definecolor{codebg}{HTML}{F5F7FA}
\definecolor{codeframe}{HTML}{D9E2EC}

\lstdefinestyle{cacheircode}{
  basicstyle=\ttfamily\small,
  backgroundcolor=\color{codebg},
  frame=single,
  rulecolor=\color{codeframe},
  breaklines=true,
  columns=fullflexible,
  keepspaces=true,
  showstringspaces=false
}

\lstset{style=cacheircode}

\title{\textbf{CacheIR Project Documentation}\\
\large A Narrow, Explainable Compiler and Runtime for Transformer Inference}
\author{CacheIR Contributors}
\date{\today}

\begin{document}
\maketitle

\begin{abstract}
CacheIR is a focused machine-learning compiler and runtime for decoder-only
transformer inference. It imports Llama, Mistral, and Qwen-style model metadata,
lowers the model into a custom intermediate representation, applies
transformer-specific optimization passes, specializes prefill and decode graphs,
plans memory and execution schedules, selects backend kernels, and executes
inference through an independent reference runtime. The project is intentionally
not a wrapper around PyTorch and not a clone of vLLM. Its primary design goal is
to make the complete path from model graph to runtime schedule inspectable,
testable, and extensible.
\end{abstract}

\tableofcontents
\newpage

\section{Project Purpose}

General ML compiler frameworks such as TVM, MLIR-based stacks, IREE, and
StableHLO-oriented pipelines are broad systems designed to support many models
and many hardware targets. Large language model serving has a narrower and very
distinct execution profile. Prefill processes an entire prompt and tends to be
compute heavy. Decode generates one token at a time and tends to be dominated by
KV-cache traffic, memory bandwidth, batching policy, and attention kernel choice.

CacheIR exists to explore that narrower problem directly. The project focuses on
five ideas:

\begin{enumerate}[leftmargin=*]
  \item Treat prefill and decode as separate compiler modes.
  \item Compile transformer operations and KV-cache policy together.
  \item Expose each optimization pass as a textual graph diff.
  \item Emit memory plans, hardware hints, and execution schedules as first-class
        artifact data.
  \item Provide a real CPU reference runtime path that does not execute through
        PyTorch, plus an optional CUDA backend where CacheIR still owns graph
        execution and PyTorch supplies the tensor/cuBLAS substrate.
\end{enumerate}

The current implementation is a complete reference project. It includes an
installable Python package, command-line interface, model importer skeletons,
CacheIR graph IR, compiler passes, artifact bundles, graph visualization,
benchmarks, a NumPy CPU runtime, an end-to-end CUDA artifact runtime, a paged
KV-cache reference implementation, an OpenAI-compatible local server with health
and metrics endpoints, prefix-cache reuse with counters, a small continuous-batch
scheduler with CUDA batched prefill/decode paths, calibrated spillover
experiments, native GGUF dense and classic
quantized block reads, StableHLO and MLIR-style round-trip experiments, a
C++20/OpenMP AVX backend with an optional pybind11 bridge, guarded Triton
kernels, optional CUDA fused-kernel sources, FP16 Tensor Core matmul experiments,
IREE/TVM upstream benchmark bridges, accelerator adapter probes, Docker support,
CI, and tests.

\section{High-Level Architecture}

CacheIR is organized into five layers:

\begin{longtable}{p{0.18\linewidth}p{0.76\linewidth}}
\toprule
\textbf{Layer} & \textbf{Role} \\
\midrule
Model importer & Reads model configuration, tensor metadata, tokenizer metadata,
rope settings, attention layout, shape information, and quantization hints. \\
Intermediate representation & Represents the model as CacheIR graph nodes with
explicit tensor types, weights, attributes, graph modes, layouts, schedules, and
memory plans. \\
Optimization passes & Transform the logical graph into optimized, mode-specific
graphs for prefill and decode. \\
Kernel backend & Selects CPU or GPU-oriented kernels and records hardware-aware
execution hints. \\
Runtime & Loads compiled artifacts, weights, tokenizers, and KV-cache state, then
executes scheduled graph nodes through backend dispatch. \\
\bottomrule
\end{longtable}

\subsection{Repository Layout}

\begin{lstlisting}
cacheir/
  ir.py                  CacheIR graph IR
  compiler.py            importer selection and pass pipeline
  importers/             HF, ONNX, GGUF, StableHLO, tiny model support
  passes/                compiler optimization passes
  runtime/               artifacts, tokenizer, KV cache, CPU/CUDA runtimes, scheduler, server
  backends/              registry, native bridge, Triton kernels, upstream probes
  benchmark.py           benchmark runner
  visualize.py           HTML/DOT/IR/MLIR graph export
  hardware.py            hardware profile discovery
  quantization.py        quantization simulation utilities
cpp/                     C++20/OpenMP AVX backend and pybind11 module
docs/                    Markdown and LaTeX documentation
examples/                tiny model compile demo
scripts/                 benchmark matrix, scheduler, GPU, and external comparison utilities
tests/                   pytest coverage
\end{lstlisting}

\section{Model Import}

\subsection{Supported Import Paths}

CacheIR currently supports several import paths at different maturity levels:

\begin{longtable}{p{0.24\linewidth}p{0.18\linewidth}p{0.50\linewidth}}
\toprule
\textbf{Format} & \textbf{Status} & \textbf{Details} \\
\midrule
Hugging Face directory & Reference path & Reads \texttt{config.json}, discovers
NPZ or safetensors tensor shapes, and builds a Llama/Mistral/Qwen-style decoder
graph. \\
NPZ tiny model & Reference test path & Used by \texttt{cacheir make-tiny}, tests,
examples, and benchmarks. \\
safetensors & Metadata path & Shape discovery and runtime tensor loading
are supported when the optional \texttt{safetensors} dependency is installed. \\
ONNX & Generic graph skeleton & Imports generic ONNX nodes into CacheIR ops where
straightforward. Full transformer semantic recovery is not claimed. \\
GGUF & Metadata plus native subset & Parses GGUF header metadata and tensor
directory, maps common Llama tensor names, builds a CacheIR graph, and can read
dense \texttt{F32}/\texttt{F16}/\texttt{BF16}/\texttt{I8}/\texttt{I16}/\texttt{I32}/\texttt{I64}/\texttt{F64} plus \texttt{Q4\_0}/\texttt{Q4\_1}/\texttt{Q5\_0}/\texttt{Q5\_1}/\texttt{Q8\_0}/\texttt{Q8\_1}
tensor payloads directly for native execution experiments. K/IQ/TQ/NV/MX
families are delegated to the optional upstream \texttt{gguf} reference package
when it exposes a dequantizer. \\
StableHLO text & Experimental subset & Parses textual StableHLO arithmetic,
shape, broadcast, slice, reduce, and cast operations and maps them into CacheIR
nodes for importer experiments. \\
\bottomrule
\end{longtable}

\subsection{Model Configuration}

The central importer data structure is \texttt{ModelConfig}. It records:

\begin{itemize}[leftmargin=*]
  \item architecture name
  \item hidden size
  \item intermediate MLP size
  \item number of layers
  \item number of attention heads
  \item number of key/value heads
  \item vocabulary size
  \item maximum position embeddings
  \item RoPE theta
  \item RMSNorm epsilon
  \item dtype
  \item tied embedding policy
  \item quantization hints
\end{itemize}

This configuration is used to construct a logical decoder graph. For each layer,
the importer emits attention norm, Q/K/V projections, RoPE, grouped-query
attention, output projection, residual adds, MLP norm, gate/up projections, SiLU,
elementwise multiplication, down projection, and final normalization/head nodes.

\section{Intermediate Representation}

\subsection{Design Goals}

CacheIR's IR is intentionally small and text-friendly. It is not trying to be a
general replacement for MLIR, StableHLO, or TVM Relay. Instead, it represents the
subset of operations needed to compile and run decoder-only transformer inference
while preserving enough metadata for scheduling, memory planning, and pass-level
explainability.

\subsection{Core Data Types}

\begin{longtable}{p{0.24\linewidth}p{0.68\linewidth}}
\toprule
\textbf{Type} & \textbf{Purpose} \\
\midrule
\texttt{TensorType} & Shape, dtype, layout, byte-size estimation, and symbolic
dimension handling. \\
\texttt{WeightSpec} & Maps an IR value to a named tensor key, tensor type,
optional quantization metadata, and optional backing file. \\
\texttt{Node} & Operation, inputs, outputs, attributes, name, and side-effect
flag. \\
\texttt{Graph} & Complete IR graph with inputs, outputs, nodes, values, weights,
constants, graph attributes, mode, and target. \\
\texttt{CompileArtifact} & Serializable package containing graphs, pass traces,
model config, target, quantization, and metadata. \\
\bottomrule
\end{longtable}

\subsection{Important IR Operations}

\begin{longtable}{p{0.30\linewidth}p{0.62\linewidth}}
\toprule
\textbf{Operation} & \textbf{Meaning} \\
\midrule
\texttt{token\_embedding} & Gathers hidden vectors from token IDs and embedding
weights. \\
\texttt{rms\_norm} & RMS normalization with learned scale. \\
\texttt{matmul} & Linear projection using out-by-in weight layout. \\
\texttt{qkv\_projection} & Fused logical Q/K/V projection produced by compiler
passes. \\
\texttt{rope} & Rotary position embedding applied to Q and K. \\
\texttt{grouped\_query\_attention} & Logical GQA/MQA attention before mode
specialization. \\
\texttt{paged\_attention\_prefill} & Prefill-specialized attention op. \\
\texttt{paged\_attention\_decode} & Decode-specialized attention op. \\
\texttt{fused\_rmsnorm\_qkv\_rope} & Fused norm, Q/K/V projection, and RoPE op. \\
\texttt{fused\_swiglu} & Fused gate projection, up projection, SiLU, and multiply. \\
\texttt{quantized\_matmul} & Quantization-aware matmul lowering. \\
\bottomrule
\end{longtable}

\subsection{Example IR}

Logical attention begins as separate nodes:

\begin{lstlisting}
%attn_norm = rms_norm(%hidden, %attn_norm_w)
%q = matmul(%attn_norm, %q_proj_w)
%k = matmul(%attn_norm, %k_proj_w)
%v = matmul(%attn_norm, %v_proj_w)
%q_rope, %k_rope = rope(%q, %k)
%attn = grouped_query_attention(%q_rope, %k_rope, %v)
\end{lstlisting}

After fusion and decode specialization, the graph contains fewer scheduled calls:

\begin{lstlisting}
%q_rope, %k_rope, %v = fused_rmsnorm_qkv_rope(...)
%attn = paged_attention_decode(%q_rope, %k_rope, %v)
\end{lstlisting}

\section{Compiler Pipeline}

\subsection{Pass Manager}

The pass manager applies a fixed transformer-oriented pipeline to each requested
execution mode. A compile call can request both \texttt{prefill} and
\texttt{decode}; each mode receives a cloned graph and its own pass trace.

\subsection{Passes}

\begin{longtable}{p{0.28\linewidth}p{0.64\linewidth}}
\toprule
\textbf{Pass} & \textbf{Effect} \\
\midrule
Shape inference & Infers tensor types for values produced by each op. \\
Constant folding & Evaluates small constant arithmetic and activation nodes. \\
QKV projection fusion & Replaces separate Q, K, and V matmuls with one logical
\texttt{qkv\_projection}. \\
RMSNorm/QKV/RoPE fusion & Fuses RMSNorm, QKV projection, and RoPE into
\texttt{fused\_rmsnorm\_qkv\_rope}. \\
SwiGLU fusion & Fuses gate/up projection, SiLU, and elementwise multiply into
\texttt{fused\_swiglu}. \\
Prefill/decode specialization & Replaces logical attention with
\texttt{paged\_attention\_prefill} or \texttt{paged\_attention\_decode}. \\
Layout conversion & Selects CPU or GPU-oriented activation and weight layouts. \\
Quantization-aware lowering & Rewrites matmul-family ops to quantized variants
when a quantization mode is requested. \\
Dead node elimination & Removes values that no longer contribute to graph
outputs. \\
Hardware-adaptive planning & Records page sizes, low-VRAM hints, and cost-class
metadata based on local hardware profile and target. \\
Kernel selection & Assigns backend kernel names to graph nodes. \\
Execution planning & Emits a schedule with step number, op, kernel, inputs,
outputs, estimated bytes, and estimated FLOPs. \\
Static memory planning & Computes a reusable activation arena with offsets,
sizes, and value lifetimes. \\
\bottomrule
\end{longtable}

\subsection{Pass Trace and Explainability}

Every pass records:

\begin{itemize}[leftmargin=*]
  \item pass name
  \item whether the graph changed
  \item remarks
  \item full before text
  \item full after text
  \item unified diff
\end{itemize}

This makes the compiler inspectable from the command line:

\begin{lstlisting}
cacheir inspect examples/tiny_artifact \
  --mode decode \
  --pass-name prefill_decode_specialization
\end{lstlisting}

\section{Prefill and Decode Specialization}

Prefill and decode are not treated as one static graph. CacheIR compiles separate
graphs because the phases stress different resources.

\subsection{Prefill}

Prefill consumes the prompt sequence. It uses causal masking across the prompt and
is relatively compute heavy. The prefill graph records
\texttt{paged\_attention\_prefill} nodes and a memory plan sized for prompt-length
activations.

\subsection{Decode}

Decode usually consumes one token per step. It appends keys and values to the
KV-cache and repeatedly attends over the cache. The decode graph records
\texttt{paged\_attention\_decode} nodes, smaller activation memory requirements,
decode-oriented page-size hints, and cost-class labels such as
\texttt{kv\_bandwidth\_bound}.

\section{KV-Cache-Aware Runtime}

\subsection{Paged KV Cache}

The runtime includes a \texttt{PagedKVCache}. The reference CPU implementation
stores keys and values in contiguous NumPy arrays for simplicity, but it also
maintains page metadata:

\begin{itemize}[leftmargin=*]
  \item page size
  \item page IDs
  \item layer-to-page mapping
  \item token offsets and page lengths
\end{itemize}

This gives the compiler and runtime a stable policy contract that Triton and CUDA
kernels can consume.

\subsection{Prefix Cache and Continuous Batching}

The runtime includes a small LRU \texttt{PrefixCache}. It stores KV-cache
snapshots keyed by token prefixes and records hit, miss, and eviction counters.
Prefill can restore a cached prefix and append only the uncached suffix while
preserving RoPE offsets and causal masking. This makes prefix reuse observable
and testable rather than a hidden serving optimization.

The continuous-batch scheduler creates independent per-request KV-cache sessions
through the runtime fork API while sharing weights, tokenizer state, and the
prefix cache. It advances active requests in decode rounds and records submitted
requests, completed requests, generated tokens, prefill/decode milliseconds,
reused prefix tokens, scheduling rounds, and maximum observed batch size. The
reference CPU backend still executes each request independently, but the
scheduling contract mirrors the control surface used by the shared GPU paged-KV
runtime.

For CUDA sessions, the scheduler can batch variable-length prompt prefill when
prefix reuse is disabled. Dense graph work runs over a padded batch, while
attention and KV writes are sliced back to each request's real sequence length.
Active decode rounds advance through \texttt{run\_decode\_batch}. The current
implementation batches norm, matmul, residual, MLP, and decode attention work.
Forked CUDA sessions share a persistent page-backed CUDA KV pool, so fused
Triton decode kernels consume a shared page table instead of repacking K/V
tensors per request. The scheduler also supports queue limits, blocking
backpressure attempts, request priorities, fairness aging, resumable preemption,
queued cancellation, request token limits, queue-wait accounting, and bounded
latency counters.

\subsection{Calibrated Spillover Cost Model}

The spillover policy can be driven by a \texttt{SpilloverCostModel}. It records
page bytes, free GPU memory, safety margin, PCIe bandwidth, CPU/GPU bandwidth
estimates, and transfer latencies. From those inputs it computes a resident-page
budget and annotates spilled pages with target and estimated transfer cost. This
keeps low-VRAM experiments deterministic in tests while still allowing hardware
profiles or measured bandwidths to tighten the policy on real machines.

\subsection{Runtime Flow}

The runtime performs the following steps:

\begin{enumerate}[leftmargin=*]
  \item Load a \texttt{CompileArtifact}.
  \item Create a \texttt{WeightStore}.
  \item Load tokenizer metadata through \texttt{TokenizerBridge}.
  \item Allocate a \texttt{PagedKVCache}.
  \item Convert input text to token IDs.
  \item Execute the prefill graph.
  \item Repeatedly execute the decode graph for greedy generation.
  \item Optionally schedule multiple requests through
        \texttt{ContinuousBatchScheduler}.
\end{enumerate}

\section{Backend Model}

\subsection{CPU Reference Backend}

The CPU backend is a NumPy reference runtime. It implements:

\begin{itemize}[leftmargin=*]
  \item token embedding
  \item RMSNorm
  \item matmul
  \item fused RMSNorm/QKV/RoPE
  \item RoPE
  \item paged attention prefill and decode
  \item residual add
  \item SiLU
  \item elementwise multiply
  \item fused SwiGLU
  \item softmax
\end{itemize}

This backend is designed for correctness, compiler validation, artifact testing,
and benchmark regression tracking. It is not presented as a high-performance LLM
serving backend.

\subsection{C++ Backend and pybind11 Bridge}

The \texttt{cpp/} directory contains a C++20/OpenMP backend with kernels for
RMSNorm, SiLU, fused SiLU-multiply, and out-by-in matmul. The matmul and RMSNorm
paths use a runtime dot-product dispatcher with scalar, AVX2/FMA, and AVX512
implementations where the compiler and host CPU support them. The static library
builds with CMake and Ninja:

\begin{lstlisting}
cmake -S cpp -B cpp/build -DCMAKE_BUILD_TYPE=Release
cmake --build cpp/build --config Release
\end{lstlisting}

The optional pybind11 module exposes native RMSNorm, fused SiLU-multiply,
out-by-in matmul, and SIMD capability reporting to Python:

\begin{lstlisting}
python -m pip install -e ".[native]"
cmake -S cpp -B cpp/build -DCACHEIR_BUILD_PYTHON=ON -DCMAKE_BUILD_TYPE=Release
cmake --build cpp/build --config Release
\end{lstlisting}

On Windows, pinning \texttt{Python3\_EXECUTABLE} to the active interpreter avoids
CMake choosing another Python installation. The latest local run installed
\texttt{pybind11} from PyPI and built \texttt{\_cacheir\_native} successfully.
The Python runtime dispatches native RMSNorm and fused SiLU-multiply
automatically when available. Set \texttt{CACHEIR\_NATIVE\_MATMUL} to
\texttt{force} or \texttt{auto} to opt into native matmul; NumPy's optimized
contraction is still faster for many local decode shapes.

\subsection{CUDA Artifact Runtime}

\texttt{CudaRuntime} executes the same compiled CacheIR artifacts on CUDA
tensors. It is not a replacement compiler pipeline and it does not hand execution
to an upstream serving engine. CacheIR still owns graph traversal, KV-cache
state, tokenizer integration, and kernel dispatch. PyTorch CUDA is used as the
GPU tensor substrate and cuBLAS access path, while CacheIR dispatches its own
Triton elementwise kernels where useful.

The CUDA runtime includes:

\begin{itemize}[leftmargin=*]
  \item lazy CPU-to-GPU weight loading with fp16 or fp32 runtime dtype selection
  \item GPU-resident KV-cache state with the same page metadata contract as the
        CPU reference cache
  \item artifact-driven prefill and decode graph execution
  \item SDPA attention dispatch, which can use optimized CUDA attention kernels
        available in the installed PyTorch build
  \item guarded Triton RMSNorm and SiLU-multiply dispatch for larger tensors
  \item cached packed QKV and Gate/Up weights, reducing decode matmul launches
        for fused CacheIR nodes
  \item persistent page-backed CUDA KV pools consumed directly by single-request
        and multi-request Triton paged decode kernels
  \item packed int4/int8 weight objects with compressed \texttt{uint8} storage,
        per-row scales, affine zero points, and CPU/CUDA loader integration
  \item shape-specific GEMM plan recording, default Torch matmul dispatch, and
        guarded opt-in Triton Tensor Core matmul
  \item CUDA memory-limit checks and OOM cleanup hooks
  \item \texttt{cacheir benchmark --backend cuda --cuda-dtype float16} support
\end{itemize}

This backend moves CacheIR beyond isolated GPU kernel microbenchmarks: complete
decoder artifacts can now execute on the RTX 5070 Laptop GPU, including a real
Qwen2.5-0.5B-Instruct CUDA prefill/decode smoke through persistent paged decode
attention.

\subsection{GPU and Triton Target Surface}

The compiler can target \texttt{cuda} or \texttt{triton} for kernel naming and
schedule generation. The project includes guarded Triton kernels for RMSNorm and
SiLU/SwiGLU activation multiply, fused RMSNorm/QKV/RoPE projection,
single-query decode attention, and multi-batch page-table decode attention with
GQA mapping when the optional \texttt{triton} package and a CUDA-capable
environment are available. It also includes a guarded Triton FP16 matmul kernel
that lowers through \texttt{tl.dot}. The \texttt{cpp/} tree includes CUDA C++
fused kernels for RMSNorm, SwiGLU, RMSNorm/QKV/RoPE, FP16 WMMA Tensor Core
matmul, a batched paged-attention ABI, a reduced block-level paged-attention
decode path, and CUDA graph capture planning.

\subsection{Accelerator Adapter Probes}

CacheIR includes guarded adapter probes and direct wrappers for CUTLASS,
FlashAttention, and FlashInfer. These adapters do not replace the compiler
pipeline: CacheIR still owns IR, memory planning, schedules, and dispatch
decisions. When an optional package is importable, a compatible scheduled op can
be routed to that external library for experiments and benchmark comparison.
CacheIR exposes guarded direct wrappers for FlashAttention prefill, FlashInfer
single decode, and FlashInfer batch paged decode wrapper classes. On the current
Windows verification machine, \texttt{nvidia-cutlass 4.2.0.0} installs and is
detected through \texttt{cutlass\_cppgen}. Native Windows/Python 3.13 still lacks
a complete wheel set for vLLM, FlashAttention, FlashInfer, and llama-cpp-python,
so the CUDA validation path uses WSL2 environments. The current WSL2 environment
executed the FlashInfer adapter smoke path; FlashAttention remains guarded
because no compatible binary wheel was available in that environment.

\section{Quantization}

CacheIR supports quantization-aware lowering. When a compile call requests a
quantization mode such as \texttt{int4\_awq}, matmul-family ops are rewritten to
quantized variants:

\begin{itemize}[leftmargin=*]
  \item \texttt{matmul} becomes \texttt{quantized\_matmul}
  \item \texttt{fused\_rmsnorm\_qkv\_rope} becomes
        \texttt{quantized\_fused\_rmsnorm\_qkv\_rope}
  \item \texttt{fused\_swiglu} becomes \texttt{quantized\_fused\_swiglu}
\end{itemize}

The CPU reference runtime simulates quantized weights by quantizing and
dequantizing weights once on load. This validates graph lowering and runtime
compatibility but is not a compressed high-performance kernel implementation.

\section{Artifacts}

\subsection{Single-File Artifact}

A \texttt{CompileArtifact} can be saved as JSON. It includes:

\begin{itemize}[leftmargin=*]
  \item artifact version and creation time
  \item target
  \item quantization mode
  \item model path
  \item model configuration
  \item graphs by mode
  \item pass traces by mode
  \item compiler metadata
\end{itemize}

\subsection{Bundle Artifact}

When the output path has no file suffix, CacheIR writes a bundle:

\begin{lstlisting}
artifact.json
manifest.json
README.txt
graphs/prefill.cir
graphs/decode.cir
schedules/prefill.json
schedules/decode.json
passes/*.diff
\end{lstlisting}

This makes the compiler output reviewable without writing custom tooling.

\section{Command-Line Interface}

CacheIR exposes the following commands:

\begin{longtable}{p{0.20\linewidth}p{0.70\linewidth}}
\toprule
\textbf{Command} & \textbf{Purpose} \\
\midrule
\texttt{profile} & Print local CPU/GPU hardware profile. \\
\texttt{make-tiny} & Generate a small Llama-shaped test model. \\
\texttt{compile} & Compile a model into a CacheIR artifact. \\
\texttt{inspect} & Print optimized IR or a specific pass diff. \\
\texttt{export} & Export IR as HTML, DOT, or text. \\
\texttt{benchmark} & Measure prefill and decode performance separately. \\
\texttt{external} & Probe optional upstream systems and run installed IREE/TVM
smoke benchmarks. \\
\texttt{run} & Run greedy local generation. \\
\texttt{serve} & Start an OpenAI-compatible local server. \\
\bottomrule
\end{longtable}

\subsection{Example Session}

\begin{lstlisting}
cacheir make-tiny examples/tiny_model
cacheir compile examples/tiny_model --output examples/tiny_artifact
cacheir inspect examples/tiny_artifact --mode decode
cacheir benchmark examples/tiny_artifact --decode-tokens 32 --repeats 3
cacheir benchmark examples/tiny_artifact --backend cuda --cuda-dtype float16 --warmup 2
cacheir external --benchmark --workdir .tmp/upstream
cacheir export examples/tiny_artifact examples/decode.html --mode decode
cacheir run examples/tiny_artifact --prompt "CacheIR"
\end{lstlisting}

\section{Python API}

The public API is intentionally small:

\begin{lstlisting}[language=Python]
from cacheir import Runtime, compile_model

artifact = compile_model(
    model_path="examples/tiny_model",
    target="cpu",
    quant=None,
    mode=["prefill", "decode"],
    max_seq=128,
)

rt = Runtime(artifact)
for token in rt.generate("Explain MLIR in simple terms", max_new_tokens=16):
    print(token, end="")
\end{lstlisting}

\section{Serving API}

When optional server dependencies are installed, CacheIR can start a FastAPI
server with OpenAI-compatible endpoints:

\begin{lstlisting}
cacheir serve examples/tiny_artifact --host 127.0.0.1 --port 8000
\end{lstlisting}

Supported endpoints include:

\begin{itemize}[leftmargin=*]
  \item \texttt{GET /v1/models}
  \item \texttt{POST /v1/completions}
  \item \texttt{POST /v1/chat/completions}
\end{itemize}

Streaming chat completions return server-sent events with OpenAI-style JSON
chunks.

\section{Benchmarking}

\subsection{Single Artifact Benchmark}

\begin{lstlisting}
cacheir benchmark examples/tiny_artifact \
  --prompt "CacheIR benchmark" \
  --decode-tokens 64 \
  --repeats 5
\end{lstlisting}

The output reports:

\begin{itemize}[leftmargin=*]
  \item prompt token count
  \item decode token count
  \item average prefill latency
  \item average decode latency per token
  \item prefill tokens per second
  \item decode tokens per second
  \item KV-cache page statistics
\end{itemize}

\subsection{CUDA Runtime Benchmark}

The dedicated CUDA runtime benchmark compares the CPU reference backend and the
CUDA artifact backend through the same benchmark API:

\begin{lstlisting}
python scripts/benchmark_cuda_runtime.py \
  --output .tmp/reports/cuda_runtime_benchmark_h1024_latest.json \
  --repeats 3 \
  --warmup 1 \
  --decode-tokens 8 \
  --hidden-size 1024 \
  --num-layers 4 \
  --cuda-dtype float16
\end{lstlisting}

{\small
\setlength{\tabcolsep}{3pt}
\begin{longtable}{@{}p{0.19\linewidth}p{0.14\linewidth}p{0.12\linewidth}p{0.14\linewidth}p{0.12\linewidth}p{0.12\linewidth}@{}}
\toprule
\textbf{Shape} & \textbf{Backend} & \textbf{Prefill ms} & \textbf{Decode ms/tok} &
\textbf{Prefill tok/s} & \textbf{Decode tok/s} \\
\midrule
h128, 2 layers & CPU fp32 & 2.122 & 1.007 & 17,432.8 & 993.0 \\
h128, 2 layers & CUDA fp16 & 3.058 & 3.210 & 12,099.6 & 311.5 \\
h512, 2 layers & CPU fp32 & 5.828 & 2.077 & 6,348.3 & 481.4 \\
h512, 2 layers & CUDA fp16 & 3.867 & 4.030 & 9,567.5 & 248.1 \\
h1024, 4 layers & CPU fp32 & 30.984 & 12.271 & 1,194.1 & 81.5 \\
h1024, 4 layers & CUDA fp16 & 7.533 & 7.573 & 4,912.0 & 132.1 \\
h2048, 4 layers & CPU fp32 & 94.036 & 30.619 & 393.5 & 32.7 \\
h2048, 4 layers & CUDA fp16 & 6.705 & 7.333 & 5,517.9 & 136.4 \\
\bottomrule
\end{longtable}
}

The CUDA backend is slower on very small toy shapes where launch overhead
dominates. It crosses over as the model becomes compute-heavy: h1024/l4 is
4.11x faster on prefill and 1.62x faster on decode, while h2048/l4 is 14.02x
faster on prefill and 4.18x faster on decode. The result proves end-to-end GPU
model execution, not production parity with mature LLM serving engines.

\subsection{Benchmark Matrix}

The repository includes \texttt{scripts/benchmark\_matrix.py}, which generates
three reference model sizes, compiles fp32 and packed quantized variants, runs
three prompt lengths, and emits JSON plus Markdown summaries:

\begin{lstlisting}
python scripts/benchmark_matrix.py \
  --output benchmark_results.json \
  --repeats 3 \
  --decode-tokens 16
\end{lstlisting}

\subsection{Scheduler Benchmark}

The repository also includes \texttt{scripts/benchmark\_scheduler.py}. It
creates or loads a CacheIR artifact, submits prompts to
\texttt{ContinuousBatchScheduler}, and reports request throughput, generated
token throughput, prefix reuse, and scheduling metrics:

\begin{lstlisting}
python scripts/benchmark_scheduler.py \
  --output .tmp/reports/scheduler_benchmark.json \
  --repeats 3 \
  --max-new-tokens 8 \
  --max-batch-size 4
\end{lstlisting}

\subsection{CUDA Scheduler Benchmark}

The CUDA scheduler benchmark compares sequential CUDA sessions with
variable-length batched prefill and active-request batched decode:

\begin{lstlisting}
python scripts/benchmark_cuda_scheduler.py \
  --output .tmp/reports/cuda_scheduler_benchmark_latest.json \
  --repeats 3 \
  --warmup 1 \
  --max-batch-size 4 \
  --max-new-tokens 16 \
  --hidden-size 1024 \
  --num-layers 4 \
  --cuda-dtype float16
\end{lstlisting}

\begin{longtable}{p{0.24\linewidth}p{0.26\linewidth}p{0.18\linewidth}p{0.18\linewidth}}
\toprule
\textbf{Shape} & \textbf{Mode} & \textbf{Median ms} & \textbf{Gen tok/s} \\
\midrule
h1024/l4, batch 4 & Sequential & 582.177 & 109.9 \\
h1024/l4, batch 4 & Batched prefill + decode & 339.796 & 188.3 \\
h2048/l4, batch 4 & Sequential & 867.240 & 37.0 \\
h2048/l4, batch 4 & Batched prefill + decode & 684.002 & 46.8 \\
\bottomrule
\end{longtable}

The h1024/l4 run improves latency and generated-token throughput by 1.71x. The
h2048/l4 run improves by 1.27x. These runs processed 134 real prefill tokens and
6 padding tokens. The shorter h2048 decode window leaves prefill as the dominant
cost, and attention still splits per request after the dense padded graph work.

\subsection{Installed Upstream Benchmarks}

When the optional upstream packages are installed, CacheIR can compile and run
small real upstream benchmark probes without user-supplied commands:

\begin{lstlisting}
cacheir external --benchmark --workdir .tmp/upstream

python scripts/compare_external_benchmarks.py examples/tiny_artifact \
  --run-installed-smoke \
  --output benchmark_comparison.json
\end{lstlisting}

The IREE path compiles StableHLO text to a VM flatbuffer through
\texttt{iree-base-compiler} and runs \texttt{iree-benchmark-module}. The TVM
path builds and runs a TE/TIR vector-add module through the installed
\texttt{apache-tvm} wheel. When model artifacts are available, CacheIR can also
launch model-aware vLLM and llama.cpp helpers:

\begin{lstlisting}
cacheir external --benchmark --workdir .tmp/upstream \
  --vllm-model /path/to/hf_model \
  --llama-model /path/to/model.gguf
\end{lstlisting}

These helpers run the installed upstream tools rather than wrapping them inside
CacheIR's runtime. The vLLM helper drives \texttt{vllm bench latency}; the
llama.cpp helper drives \texttt{llama-bench}.

\subsection{Local Native Experiment Snapshot}

The latest local verification on 2026-06-29 exercised the native, CUDA, WSL2,
and upstream comparison paths. The full test suite passed with 38 tests, and
\texttt{python -m compileall cacheir scripts -q} completed after the production
runtime changes.

The native pybind11 module built successfully for Python 3.13 and reported
\texttt{simd\_backend() == "avx512"}. Native RMSNorm, matmul, and fused
SiLU-multiply matched NumPy within float32 tolerance. Packed int4 and int8
weight objects now store compressed uint8 payloads, per-row scales, affine zero
points, and original shape metadata; CPU and CUDA loaders route quantized
matmul, QKV, and SwiGLU paths through those objects.

The CUDA runtime now maintains persistent page-backed GPU KV pools. The fused
Triton decode path consumes the shared physical page table directly for
multi-request decode rather than repacking per-request K/V tensors. On h512/l2
with \texttt{--use-triton-attention}, scheduler median latency improved from
176.225 ms to 78.491 ms and generated-token throughput improved from
184.0 tok/s to 407.8 tok/s. A h512/l2 end-to-end CUDA runtime benchmark
measured 3.515 ms prefill versus CPU at 8.800 ms; decode was 3.264 ms/token on
CUDA versus 3.139 ms/token on CPU for that small shape.

The Triton kernel benchmark on the RTX 5070 Laptop GPU measured RMSNorm at
0.0185 ms, SiLU-multiply at 0.0189 ms, FP16 matmul at 0.0674 ms or
7.97 TFLOP/s, Torch FP16 matmul at 0.0281 ms or 19.14 TFLOP/s, and persistent
multi-batch page-table decode attention at 0.0583 ms. The scheduler hardening
pass added blocking backpressure, fairness aging, preemption, per-request token
limits, bounded latency counters, and long-running queue soak metrics.

Qwen2.5-0.5B-Instruct was downloaded locally in WSL2, compiled through the
Hugging Face importer, and executed on CUDA with bf16 safetensors shape
discovery and persistent paged decode. For a 16 input + 8 output token request,
CacheIR averaged 233.774 ms/request after warmup, while vLLM 0.23.0 with Torch
2.11.0+cu130 and FlashInfer averaged 43.697 ms/request in the same WSL2 CUDA
environment. vLLM remains substantially faster on this model because it uses
mature CUDA graphs, lower Python overhead, and heavily tuned kernels.

A CUDA llama.cpp build ran \texttt{llama-bench} against a locally converted tiny
GGUF Llama model: 6,120.22 prompt tok/s for 16 prompt tokens and 2,044.44
generation tok/s for 8 generated tokens. \texttt{iree-base-compiler 3.11.0},
\texttt{iree-base-runtime 3.11.0}, and \texttt{apache-tvm 0.25.0} installed from
PyPI; the upstream benchmark command compiled StableHLO through IREE to a
9,781-byte VMFB and ran a TVM TE vector-add benchmark with checksum 512.0.
FlashInfer executed through CacheIR's adapter smoke test. FlashAttention,
TensorRT-LLM, and MLC LLM are guarded optional integrations on this host because
binary-only PyPI checks found no compatible wheels and source install attempts
were not stable in the local WSL2 image.

\section{Testing and Validation}

The test suite covers:

\begin{itemize}[leftmargin=*]
  \item decoder graph compilation
  \item decode specialization
  \item pass fusion behavior
  \item artifact JSON round trips
  \item artifact bundle loading
  \item graph export
  \item runtime prefill and decode
  \item CUDA runtime prefill/decode parity with the CPU reference when CUDA is
        available
  \item CUDA scheduler variable-length batched prefill, persistent shared page
        allocator accounting, and active-request batched decode parity with
        sequential CUDA sessions
  \item persistent Triton multi-request decode over the shared CUDA page pool
        when CUDA and Triton are available
  \item scheduler queue limits, request priorities, queued cancellation,
        blocking backpressure, fairness aging, preemption, and request token
        limits
  \item prefix-cache-aware prefill equivalence
  \item KV-cache growth
  \item greedy generation
  \item continuous-batch scheduler metrics
  \item server health, metrics, and batch completions
  \item benchmark execution
  \item hardware profile serialization
  \item packed int4/int8 quantized lowering, round trips, and runtime
        compatibility
  \item GGUF dense and \texttt{Q4\_0}/\texttt{Q4\_1}/\texttt{Q5\_0}/\texttt{Q5\_1}/\texttt{Q8\_0}/\texttt{Q8\_1} tensor reads
  \item StableHLO import with reduce-region metadata and MLIR-style dialect export/parser/verifier round trip
  \item calibrated spillover cost model and adapter probe behavior
  \item native pybind11 RMSNorm, matmul, and fused SiLU-multiply parity when the extension is built
\end{itemize}

The latest full local test run passed 38 tests.

Recommended validation commands:

\begin{lstlisting}
python -m pytest -q
python -m compileall cacheir scripts -q
cmake -S cpp -B cpp/build -DCMAKE_BUILD_TYPE=Release
cmake --build cpp/build --config Release
\end{lstlisting}

\section{Complete Technology Stack}

\begin{longtable}{p{0.26\linewidth}p{0.64\linewidth}}
\toprule
\textbf{Area} & \textbf{Technologies} \\
\midrule
Primary languages & Python 3.10+, C++20. \\
Numerical runtime & NumPy reference kernels plus optional PyTorch CUDA tensor
substrate for the CacheIR CUDA executor. \\
Compiler IR & Custom CacheIR graph IR with JSON and textual representations. \\
Model import & Hugging Face config conventions, NPZ, safetensors optional,
ONNX optional, GGUF metadata plus dense \texttt{F32}/\texttt{F16}/\texttt{BF16}/\texttt{I8}/\texttt{I16}/\texttt{I32}/\texttt{I64}/\texttt{F64} and
\texttt{Q4\_0}/\texttt{Q4\_1}/\texttt{Q5\_0}/\texttt{Q5\_1}/\texttt{Q8\_0}/\texttt{Q8\_1} tensor subset, StableHLO textual subset. \\
Transformer semantics & Llama/Mistral/Qwen-style decoder-only blocks,
RMSNorm, RoPE, grouped-query attention, SwiGLU, KV cache. \\
Optimization passes & Shape inference, constant folding, operator fusion,
layout conversion, dead node elimination, mode specialization, quantization-aware
lowering, schedule and memory planning. \\
CPU backend & NumPy reference runtime, C++20/OpenMP native library, scalar,
AVX2/FMA, and AVX512 dispatch, optional pybind11 bridge, native RMSNorm and
fused SiLU-multiply dispatch, opt-in native matmul policy. \\
GPU backend surface & CUDA/Triton target naming, schedule metadata,
\texttt{CudaRuntime}, fp16 CUDA weight loading, GPU-resident KV state, SDPA
attention dispatch, cached packed QKV/Gate-Up weights, shared CUDA page
allocator accounting, variable-length CUDA prefill batching,
scheduler-integrated CUDA decode batching, guarded Triton RMSNorm,
SiLU/SwiGLU, FP16 matmul, fused RMSNorm/QKV/RoPE, single-query decode attention,
and multi-batch page-table decode-attention kernels, optional CUDA C++ fused
kernels, FP16 WMMA Tensor Core matmul, reduced paged-attention decode, batched
paged-attention ABI, and CUDA graph planning. \\
Accelerator adapters & Optional CUTLASS, FlashAttention, and FlashInfer probes,
dispatch contracts, and guarded direct wrappers. \\
Quantization & int4/int8-style lowering and CPU-side quantize/dequantize
simulation. \\
Serving & FastAPI and Uvicorn optional dependencies, OpenAI-compatible
completions/chat, streaming, \texttt{/healthz}, \texttt{/metrics}, CacheIR batch
completions, queue limits, priorities, cancellation, and admission counters. \\
Testing & pytest, compileall, CMake build checks. \\
Benchmarking & CacheIR benchmark CLI with CPU/CUDA backend selector, benchmark
matrix script, CUDA runtime benchmark, CUDA scheduler benchmark, scheduler
benchmark, GPU kernel benchmark, and comparison harness for vLLM, llama.cpp,
IREE, and TVM commands plus installed IREE/TVM smoke benchmark execution and
model-aware vLLM/llama.cpp helpers. \\
Visualization & HTML, DOT, text IR, and MLIR-style CacheIR dialect export plus
parser/verifier round trip. \\
Build and packaging & pyproject.toml, setuptools, CMake, Ninja, Docker,
GitHub Actions. \\
External context & vLLM, llama.cpp, IREE, TVM, StableHLO, MLIR, CUTLASS,
FlashAttention, and FlashInfer are comparison or optional integration surfaces,
not default runtime dependencies. \\
\bottomrule
\end{longtable}

\section{Limitations and Roadmap}

\subsection{Implemented and Tested}

\begin{itemize}[leftmargin=*]
  \item CPU reference execution without PyTorch.
  \item End-to-end CUDA artifact execution through \texttt{CudaRuntime}.
  \item CUDA fp16 weight loading, GPU-resident KV state, SDPA attention dispatch,
        and cached packed QKV/Gate-Up weights.
  \item Scheduler-integrated variable-length CUDA prefill batching and
        active-request CUDA decode batching.
  \item Shared CUDA page allocator accounting across forked request sessions.
  \item Prefill and decode graph specialization.
  \item Paged KV-cache reference behavior.
  \item Prefix-cache reuse with counters and calibrated spillover policy experiments.
  \item Forked per-request runtime sessions sharing weights and tokenizer state.
  \item Continuous-batch scheduler with prefix reuse metrics, queue limits,
        priorities, cancellation, and admission-control counters.
  \item Server health, Prometheus-style metrics, streaming, and batch completions.
  \item Artifact bundles, pass diffs, graph export, benchmarks, and server
        entrypoint.
  \item Quantization-aware lowering with CPU-side quantize/dequantize simulation.
  \item C++20 backend library build with scalar, AVX2/FMA, and AVX512 dispatch.
  \item Optional pybind11 bridge for native RMSNorm, matmul, and fused
        SiLU-multiply kernels.
  \item Native matmul policy with NumPy default and opt-in native force/auto modes.
  \item Guarded Triton RMSNorm, SiLU/SwiGLU, fused RMSNorm/QKV/RoPE,
        single-query decode attention, and multi-batch page-table decode
        attention kernels.
  \item Guarded Triton FP16 matmul using \texttt{tl.dot} Tensor Core lowering
        where available.
  \item Optional CUDA C++ fused kernels for RMSNorm, SwiGLU, RMSNorm/QKV/RoPE,
        FP16 WMMA Tensor Core matmul, batched paged-attention ABI, and reduced
        paged-attention decode.
  \item CUDA graph capture planning for decode replay loops.
  \item Native GGUF dense numeric tensor reads plus classic Q4/Q5/Q8 block
        tensor reads.
  \item Reference GGUF dequantization for upstream-supported K/IQ/TQ/NV/MX formats.
  \item StableHLO text importer with reduce-region metadata.
  \item MLIR-style CacheIR dialect emitter, parser, verifier, and optional
        upstream MLIR C++ dialect registration skeleton.
  \item CUTLASS, FlashAttention, and FlashInfer adapter probes plus guarded
        direct wrappers for prefill, single decode, and batch paged decode.
  \item Direct FlashInfer smoke execution on WSL2 CUDA, with FlashAttention kept
        as a guarded adapter for environments that provide a compatible wheel.
  \item vLLM no-UVA fallback, CacheIR project-root propagation into spawned
        workers, CUDA Toolkit 13.3 path export, and runnable WSL2 CUDA vLLM
        comparison.
  \item Calibrated low-VRAM CPU/GPU spillover cost model with measured bandwidth calibration.
  \item IREE StableHLO compile/runtime benchmark integration through upstream
        IREE wheels.
  \item TVM TE/TIR runtime benchmark integration through the upstream TVM wheel.
  \item Benchmark comparison harness for vLLM, llama.cpp, IREE, and TVM commands,
        including installed IREE/TVM smoke runs and model-aware vLLM/llama.cpp
        helpers.
\end{itemize}

\subsection{Production and Performance Completion Pass}

This pass completed the previously planned production-serving and raw-performance
surfaces that fit the local hardware and package environment:

\begin{itemize}[leftmargin=*]
  \item Persistent page-backed shared GPU KV storage consumed directly by
        CacheIR's Triton single-request and multi-request decode kernels.
  \item Fused multi-request paged attention over compact per-layer physical
        page slots without K/V repacking.
  \item Packed int4/int8 weights with real compression, scales, zero points,
        and CPU/CUDA model-loader integration.
  \item Shape-specific GEMM plan recording, default Torch matmul dispatch, and
        guarded Triton Tensor Core matmul.
  \item Scheduler hardening for backpressure, fairness aging, preemption,
        bounded queue-latency counters, request token limits, and cancellation.
  \item Reliability hardening for CUDA memory limits, OOM cleanup, tokenizer
        edge cases, BF16 safetensors loading, and Qwen model-family conformance.
  \item Qwen2.5-0.5B-Instruct CUDA smoke and latency measurement on the local
        RTX 5070 Laptop GPU.
  \item Real vLLM, llama.cpp, IREE, and TVM comparison runs where those systems
        are installed locally.
\end{itemize}

The local WSL2 CUDA environment has vLLM and FlashInfer, and FlashInfer executed
through CacheIR's adapter smoke. FlashAttention, TensorRT-LLM, and MLC LLM remain
guarded optional paths here because binary-only PyPI checks did not expose
compatible wheels for the active WSL venv. Qwen2.5-1.5B and 7B-class model
weights were not already local; the implemented importer/runtime path can run
them when local artifacts are present, subject to VRAM and quantization limits.

\section{Building PDF or ODF Outputs}

The source of this document is intended to be version controlled. Generated
outputs are intentionally ignored by \texttt{.gitignore}.

To build a PDF with a local LaTeX distribution:

\begin{lstlisting}
pdflatex -interaction=nonstopmode docs/latex/cacheir_project_documentation.tex
\end{lstlisting}

To build an OpenDocument text file with Pandoc:

\begin{lstlisting}
pandoc docs/latex/cacheir_project_documentation.tex \
  -o docs/latex/cacheir_project_documentation.odt
\end{lstlisting}

\section{Conclusion}

CacheIR is a compact but substantial transformer inference compiler project. It
prioritizes clarity, pass-level explainability, prefill/decode specialization,
and KV-cache-aware runtime behavior. The NumPy CPU backend establishes a correct
reference path, while the artifact format, schedules, C++ AVX backend, optional
pybind11 bridge, CUDA artifact runtime, GGUF/StableHLO/MLIR experiments, guarded
Triton kernels, CUDA fused-kernel sources, accelerator adapter probes, and
calibrated spillover model, IREE StableHLO bridge, and TVM benchmark bridge
provide concrete native performance and upstream comparison surfaces. The
project is now strong as an explainable compiler and systems artifact, but
vLLM remains much faster on the measured Qwen2.5-0.5B request. CacheIR's current
value is the transparent compiler pipeline and realistic serving surface; future
performance gains should focus on reducing Python graph overhead, fusing
quantized GEMM instead of dequantizing at the GEMM boundary, and expanding local
large-model benchmark coverage.

\end{document}
